% Architektur — claude-tunnel bachelor thesis.
\chapter{Architektur}
\label{ch:architektur}

Dieses Kapitel beschreibt die Struktur des Systems: die Komponenten und ihre
Topologie (Abschnitt~\ref{sec:arch-topologie}), die zentralen Ablaeufe
(\ref{sec:arch-fluesse}), den Rollen-Dispatch auf einem QUIC-Stream
(\ref{sec:arch-dispatch}) und die tragenden Entwurfsentscheidungen
(\ref{sec:arch-entscheidungen}). Es setzt die Anforderungen aus
Kapitel~\ref{ch:anforderungen} in eine Aufteilung von Verantwortlichkeiten um.

\section{Komponenten und Topologie}
\label{sec:arch-topologie}

Das System besteht aus vier Rollen auf dem Datenpfad -- Client, Edge, Agent und
Origin -- sowie einer Steuerebene (Control Plane). Der Client will einen lokalen
Dienst erreichen; das Origin ist der Zieldienst; der Agent ist kundenseitig und
haelt den Origin-Schluessel; der Edge ist der providerblinde Vermittler.
Abbildung~\ref{fig:topologie} zeigt die Topologie.

\begin{figure}[t]
  \centering
  \begin{tikzpicture}[
      node distance=14mm and 20mm,
      box/.style={draw, rounded corners, minimum height=9mm, minimum width=18mm,
                  align=center, font=\small},
      relay/.style={-{Stealth[]}, thick},
      e2e/.style={-{Stealth[]}, dashed},
    ]
    \node[box] (client) {Client};
    \node[box, right=of client] (edge) {Edge\\\scriptsize (Relay)};
    \node[box, right=of edge] (agent) {Agent};
    \node[box, right=of agent] (origin) {Origin};
    \node[box, above=of edge] (cp) {Control\\Plane};

    \draw[relay] (client) -- node[above,font=\scriptsize]{QUIC} (edge);
    \draw[relay] (edge) -- node[above,font=\scriptsize]{QUIC} (agent);
    \draw[relay] (agent) -- node[above,font=\scriptsize]{TCP/UDP} (origin);
    \draw[relay, dotted] (agent) -- (cp);
    \draw[relay, dotted] (client) to[bend left=12] (cp);

    % Provider-blind end-to-end Noise arc: client <-> origin
    \draw[e2e] (client) to[bend right=28]
      node[below,font=\scriptsize]{Noise-Ende-zu-Ende (nur Chiffretext am Edge)}
      (origin);
  \end{tikzpicture}
  \caption{Topologie: Der Edge relayt zwischen Client und Agent, sieht aber nur
  den Ende-zu-Ende mit Noise verschluesselten Verkehr (gestrichelter Bogen). Die
  Steuerebene koordiniert (gepunktet), traegt aber keine Nutzlast.}
  \label{fig:topologie}
\end{figure}

Die Trennung ist bewusst: der Edge kennt nur opake Routing-Tokens und relayt
Bytes, das Origin terminiert die Noise-Sitzung, und die Steuerebene haelt weder
Schluessel noch Nutzlast (thin control plane, N2).

\section{Schluesselfluesse}
\label{sec:arch-fluesse}

\paragraph{Enrollment und Anmeldung.} Der Agent redeemt ein Join-Token bei der
Steuerebene und bindet seinen ed25519-Schluessel; er erhaelt kurzlebige
Anmeldedaten, die der Edge zustandslos verifiziert (F2).

\paragraph{Rendezvous mit Proof-of-Work.} Oeffnet ein Client den Tunnel, sendet
der Edge eine \gls{ac:pow}-Challenge; der Client antwortet mit Loesung und
Routing-Token. Erst nach gueltiger Loesung loest der Edge das Token zum Agent auf
(F3, F4).

\paragraph{Relay und Ende-zu-Ende-Handshake.} Der Edge verbindet den Client-Stream
mit dem registrierten Agent und kopiert fortan nur Bytes. Darueber laeuft der
Noise-Handshake zwischen Client und Origin; der Client pinnt die
Origin-Identitaet (S1, S3).

\paragraph{Direkter Pfad und Fallbacks.} Advertisiert der Agent einen direkten
Listener, entdeckt der Client dessen Adresse und Zertifikat ueber den Edge und
verbindet direkt; scheitert das, faellt er auf den Relay zurueck. Ist
\gls{ac:udp} blockiert, traegt ein \gls{ac:tls}-ueber-\gls{ac:tcp}-Kanal denselben
Tunnel (F5).

\section{Rollen-Dispatch auf einem Stream}
\label{sec:arch-dispatch}

Client und Agent teilen sich das QUIC-Protokoll zum Edge; die gewuenschte Rolle
wird durch ein fuehrendes Byte auf dem Stream ausgewaehlt:

\begin{description}
  \item[\texttt{'A'}] Agent registriert einen Tunnel (Token + reflexive Adresse).
  \item[\texttt{'C'}] Client fordert Rendezvous an (PoW), wird geroutet und
    relayt.
  \item[\texttt{'D'}] Agent advertisiert seinen direkten Listener (Adresse,
    Zertifikat).
  \item[\texttt{'P'}] Client fragt den direkten Endpunkt eines Tokens ab.
\end{description}

Dieser minimale Dispatch haelt den Edge einfach und erweiterbar, ohne ein
separates Steuerprotokoll auf eigener Verbindung.

\section{Entwurfsentscheidungen}
\label{sec:arch-entscheidungen}

\begin{description}
  \item[Thin Control Plane.] Die Steuerebene koordiniert nur (Enrollment,
    Registry, Abrechnung) und haelt kein Vertrauensmaterial; ein Kompromiss der
    Steuerebene gefaehrdet damit nicht die Vertraulichkeit der Tunnel.
  \item[Capability ausserhalb des Bandes.] Die Verbindungserlaubnis (Routing-Token
    + Origin-Identitaet + Edge-Adresse) wird vom Kunden out-of-band verteilt;
    Besitz genuegt, um das Origin zu erreichen und zu authentisieren.
  \item[Noise\_IK mit Origin-Pinning.] Die Wahl des IK-Musters verankert die
    Origin-Authentizitaet im Client und macht den Edge als Man-in-the-Middle
    wirkungslos.
  \item[Transport-agnostischer Datenpfad.] Ein gemeinsamer, bidirektionaler
    Noise-Streaming-Baustein traegt Streaming, Datagramm, direkten Pfad und
    TCP-Fallback gleichermassen; die Betriebsarten unterscheiden sich nur an den
    Raendern.
\end{description}
